\section{Related Work}

\paragraph{State Learning}

Applying model learning to analyze the security of cryptographic protocols was first proposed by de Ruiter~et~al.~\cite{USENIX:dRuPol15} in the context of the Transport Layer Security (TLS) protocol. The authors coined the term \emph{protocol state fuzzing} for this approach, as protocol state transitions are randomly learned in black box testing, resulting in a state machine that can be analyzed manually or with model checking. Various studies further extended the analysis for TLS and Datagram TLS~\cite{SecITSys:dRu16,yadav2018tls,USENIX:FJMRSS20,CCS:STVHBC22,ESORICS:RasLevDeb22}. Karim~et~al. studied the Bluetooth Low Energy (BLE) protocol~\cite{SP:KIHB23} and proposed a methodology to learn phases of the protocol separately but merge them into one comprehensive state machine later. In 2024, Tu~et~al.~\cite{USENIX:TIRDWW24} studied the state machines of 5G basebands using a novel approach combining passive and active state learning and by leveraging the models of different basebands for their equivalence oracle. Dong~et~al. introduced the concept of guiding state machines from protocol specifications to implement a mutation-based testing strategy for 4G LTE and BLE devices in~\cite{CCS:RWTITD24}. They applied a related concept to 5G core network implementations in~\cite{dong2025corecrisis} using learned state machines, thereby deriving test cases that can evaluate protocol states in greater detail.

In 2017, Fiter\u{a}u-Bro\c{s}tean~et~al. first applied state learning to SSH~\cite{SPIN:fiterauSshModelChecking}. They extracted state machines of three implementations and applied model-checking techniques to test if each state machine achieves certain functional and security properties. Their learner commenced after the exchange of the ASCII banner; however, they further partitioned the protocol into three layers (transport layer, authentication layer, and connection layer) and learned each layer individually with a distinct alphabet designed for the respective protocol phase. In 2023, Fiter\u{a}u-Bro\c{s}tean~et~al. further expanded their approach~\cite{NDSS:FJST23} and proposed a methodology for the automated validation of state machines through manually modeled bug patterns. For the SSH part of this paper, they used the three extracted state machines from~\cite{SPIN:fiterauSshModelChecking} as artifacts.

\paragraph{SSH Security Analyses} 

Various studies have presented practical analyses on the security of the SSH protocol, its implementations, and deployments on the Internet. For instance, Albrecht~et~al.~\cite{SP:AlbPatWat09} demonstrated how an attacker can use encrypted length fields to uncover the plaintext of a message. Song~et~al.~\cite{USENIX:SonWagTia01} showed that keystroke timing attacks allow an attacker to deduce passwords typed by a user during interactive authentication in a shell. Bhargavan and Leurent~\cite{NDSS:BhaLeu16} analyzed the impact of SHA-1 collisions on the session transcript of SSH. In a study evaluating the host keys of SSH deployments on the Internet, Heninger~et~al.~\cite{USENIX:HDWH12} found that implementation errors and lack of entropy in generated key material leave some hosts vulnerable to impersonation attacks. Bäumer~et~al.~\cite{USENIX:BauBriSch24} presented the Terrapin attack, a prefix truncation attack against some SSH cipher modes. They also give a rogue session attack similar to the one in \autoref{sec:eval:sub:tectia}. The countermeasure they propose is the subject of our study. Formal proofs for SSH have been contributed by Bellare~et~al.~\cite{CCS:BelKohNam02}, Bergsma~et~al.~\cite{CCS:BDKSS14}, Paterson and Watson~\cite{EC:PatWat10}, Williams~\cite{IMA:Williams11}, and Hansen~\cite{hansen2020cryptographic}.